\chapter{Using the primitives}
\label{ch:primitives}

nornCore offers more than one way to signal and to share. This chapter says what each
primitive is for, how it behaves at the edges, and shows the idiom. The formal contracts
are in chapter~\ref{ch:api}; the costs in chapter~\ref{ch:perf}.

\begin{table}[htbp]
\centering\small
\begin{tabularx}{\linewidth}{@{}lXl@{}}
\toprule
Need & Primitive & Cost, Mega / ESP32 \\
\midrule
wake one task, no data & task notification (\code{rtos\_notify}) or event flags & 220 / 145, 217 / 140 \\
count occurrences, producer to consumer & counting semaphore, or \code{RTOS\_NOTIFY\_INC} & 145 / 112 \\
hand a small record to another task, in order & message queue & 453 / 285 \\
the latest value only, overwrite is fine & mailbox & 282 / 194 \\
flags several tasks wait on & event group & 272 / 137 \\
protect shared data between tasks & ceiling mutex & 218 / 156 \\
fixed size buffers, bounded time & memory pool & 487 / 176 \\
variable size, rarely, at start-up & heap & -- \\
run something later or periodically without a task & software timer & 241 / 124 \\
\bottomrule
\end{tabularx}
\caption{Which primitive for which need. The costs are cycles for one call pair without a
switch (tc10, \measDate).}
\label{tab:which}
\end{table}

\section{Task event flags}

Every task owns \code{RTOS\_CFG\_EVT\_BITS} private event flags. Anyone posts
(\code{rtos\_postEvents}, \code{rtos\_postEventsFromIsr}, \code{rtos\_broadcastEvents}),
only the owner waits. Events posted earlier stay pending until a wait consumes them, so
nothing is lost between two waits; the wait returns the events that satisfied it and
clears those, leaving other pending events alone.

\begin{lstlisting}
#define EVT_SAMPLE  0x01u
#define EVT_BUTTON  0x02u

static void controller(void *arg)
{
    for(;;)
    {
        const rtos_evt_t got = rtos_waitEvents(EVT_SAMPLE | EVT_BUTTON, RTOS_WAIT_ANY, RTOS_MS(500));
        if(got == 0)            /* timeout: nothing in 500 ms */
            handleSilence();
        if(got & EVT_SAMPLE)
            processSample();
        if(got & EVT_BUTTON)
            handleButton();
    }
}

RTOS_ISR(ADC_vect)
{
    rtos_postEventsFromIsr(TASK_CONTROLLER, EVT_SAMPLE);
}
\end{lstlisting}

\code{RTOS\_WAIT\_ALL} waits for every bit of the mask. Event flags are the right tool
when one task is the hub of several sources; they carry no data and do not count.

\section{Counting semaphores}

A semaphore counts tokens. \code{rtos\_semGive} adds one, or hands it directly to the
highest waiting task; \code{rtos\_semTake} takes one or waits. \code{RTOS\_SEM\_INIT(n)}
makes a semaphore without an upper limit, \code{RTOS\_SEM\_INIT\_BINARY(n)} one whose
count never exceeds one (a give while the token is available is dropped and counted in
the \code{lost} field), \code{RTOS\_SEM\_INIT\_MAX(n, max)} one with any limit. The
typical use is an interrupt producing and a task consuming:

\begin{lstlisting}
static rtos_sem_t _rxReady = RTOS_SEM_INIT(0);

RTOS_ISR(USART0_RX_vect)
{
    _rxBuf[_rxIn++ & RX_MASK] = UDR0;
    rtos_semGiveFromIsr(&_rxReady);
}

static void reader(void *arg)
{
    for(;;)
        if(rtos_semTake(&_rxReady, RTOS_FOREVER))
            consume(_rxBuf[_rxOut++ & RX_MASK]);
}
\end{lstlisting}

The semaphore here is only the count; the data travels in the application's own ring. A
semaphore is not a lock: there is no owner and no priority ceiling, and using one for
mutual exclusion brings back the priority inversion the mutex exists to prevent.

\section{Ceiling mutexes}

A mutex protects shared data between tasks. Its one parameter is the ceiling, the highest
priority of any task that will lock it (section~\ref{sec:mutex}).

\begin{lstlisting}
#define TASK_LOG     1
#define TASK_SENSOR  2
#define TASK_CTRL    4

static rtos_mutex_t _stateLock = RTOS_MUTEX_INIT(TASK_CTRL);   /* highest user */
static struct state _state;

static void sensor(void *arg)        /* priority 2 */
{
    for(;;)
    {
        (void)rtos_waitPeriodic(RTOS_MS(10));
        const int16_t t = readTemperature();
        rtos_mutexLock(&_stateLock);     /* runs at 4 until the unlock */
        _state.temperature = t;
        rtos_mutexUnlock(&_stateLock);
    }
}
\end{lstlisting}

While the sensor task holds the lock it runs at priority 4, so the controller (4) cannot
preempt it and find the state half written, and the logger (1) is unaffected. The
critical section is a few stores; keep it that way. Do not call anything that may block
inside it.

\section{Message queues}

A queue is a FIFO of fixed size items copied by value, with blocking send and receive and
non-blocking \code{FromIsr} variants. A send to a queue with a blocked receiver copies
the item straight into the receiver's buffer. A full queue makes the sender wait or, with
\code{RTOS\_NO\_WAIT} or from an interrupt, fail; failed sends are counted in the
\code{dropped} field, and the high water mark in \code{maxCount}, so a queue sized too
small is visible.

\begin{lstlisting}
typedef struct { uint8_t channel; uint16_t value; } sample_t;

static uint8_t      _qStore[8 * sizeof(sample_t)];
static rtos_queue_t _q = RTOS_QUEUE_INIT(_qStore, sizeof(sample_t), 8);

RTOS_ISR(BOARD_APP_TIMER_VECT)
{
    sample_t s = { .channel = ADMUX & 0x0F, .value = ADC };
    (void)rtos_queueSendFromIsr(&_q, &s);        /* false if full: counted as dropped */
}

static void worker(void *arg)
{
    sample_t s;
    for(;;)
        if(rtos_queueReceive(&_q, &s, RTOS_MS(100)))
            process(&s);
        else
            noSampleFor100ms();
}
\end{lstlisting}

Items are copied twice (in and out), so a queue of large items is better replaced by a
queue of pointers into a memory pool.

\section{Mailboxes}

A mailbox is a queue of depth one in which a post never blocks and overwrites an unread
item (counted in \code{overwrites}). It carries \emph{the latest value}: a setpoint, a
measurement, a state that only the newest matters of. A reader blocks until something is
posted. Use it where a queue would only fill with stale data.

\section{Task notifications}

The cheapest signal: one 16 bit value per task and no object. \code{rtos\_notify}
replaces, ORs into or increments the value and wakes the task if it waits;
\code{rtos\_notifyWait} takes the whole value, \code{rtos\_notifyTake} waits for a
non-zero value and decrements it, which with \code{RTOS\_NOTIFY\_INC} on the sending side
is a counting semaphore without a semaphore. A task should use one of the two wait styles,
since they share the value. Notifications are the fastest switch on both boards
(figure~\ref{fig:switch}).

\begin{lstlisting}
RTOS_ISR(DMA_DONE_vect)
{
    rtos_notifyFromIsr(TASK_DMA, 1, RTOS_NOTIFY_INC);
}

static void dmaTask(void *arg)
{
    for(;;)
        if(rtos_notifyTake(RTOS_FOREVER))
            handleCompletion();
}
\end{lstlisting}

\section{Software timers}

A software timer calls a function after a delay, once or periodically, without a task of
its own. With \code{RTOS\_CFG\_TIMER\_TASK\_PRIO 0} the callback runs in the tick
interrupt with interrupts disabled: keep it short and use only the ISR body functions.
With a timer task configured, the callback runs in that task and may use the whole API,
at the price of a task slot and its stack.

\begin{lstlisting}
static rtos_timer_t _blink;

static void blinkCb(void *arg)
{
    gpio_toggle(&_led);                  /* ISR body: no blocking, no switch */
}

/* in a task or in main() before rtos_start: */
rtos_timerInit(&_blink, blinkCb, NULL);
rtos_timerStart(&_blink, RTOS_MS(500), RTOS_MS(500));   /* first in 500 ms, then every 500 ms */
\end{lstlisting}

The first expiry is on the \code{delay}-th tick from now (so after $\mathit{delay}-1$ to
$\mathit{delay}$ periods); later expiries are exactly \code{period} ticks apart, counted
from the previous expiry, so a periodic timer does not drift. A delay of 0 means the next
tick, which lets a callback restart its own timer. The active timers are a list sorted
by expiry; \code{rtos\_timerStart} spends time proportional to their number with
interrupts off.

\section{Memory pools}

A pool hands out fixed size blocks from storage the application owns, in constant time,
with a wait for a free block if wanted and a \code{FromIsr} variant that does not wait.
\code{rtos\_poolFree} checks that the block belongs to the pool and lies on a block
boundary. The low water mark \code{minFree} and the count of failed allocations
\code{failed} show whether the pool is sized right. A pool is what a driver uses for
packet buffers and what a queue of pointers refers to; it is the bounded time
alternative to the heap.

\begin{lstlisting}
static uint8_t     _frames[6 * 64];
static rtos_pool_t _framePool = RTOS_POOL_INIT(_frames, 64, 6);
/* rtos_poolInit(&_framePool, _frames, 64, 6) must still be called before use */

uint8_t *p = rtos_poolAlloc(&_framePool, RTOS_MS(10));
if(p != NULL) { fill(p); rtos_queueSend(&_frameQ, &p, RTOS_FOREVER); }
\end{lstlisting}

\section{Event groups}

An event group is a set of flags shared by any number of tasks
(\code{RTOS\_CFG\_EVGROUPS}). Where task event flags belong to one task, a group lets
several tasks wait for the same conditions: ``link up and address assigned'', say.
\code{rtos\_evgroupSet} wakes every waiter whose condition is now met, highest first; a
waiter that asked for \code{RTOS\_WAIT\_CLEAR} takes its bits out of the group as it is
woken, so a lower waiter no longer sees them. Without it the bits stay set for everyone
until \code{rtos\_evgroupClear}.

\section{Suspend and resume}

\code{rtos\_suspend} keeps a task from running until \code{rtos\_resume}, whatever it was
doing. A blocked task stays in its wait; a wake-up that arrives meanwhile is kept and
takes effect at the resume. There is no nesting count, and a suspended task keeps the
mutexes it holds, which is the reason to suspend only tasks that are known to hold none.
Suspension is for modes (``the display task is off while the lid is closed''), not for
synchronisation; the primitives above are.

\section{The heap}

\code{RTOS\_CFG\_HEAP} compiles a first fit allocator with coalescing over an area the
application gives it (\code{rtos\_heapInit}). \code{rtos\_malloc} searches the free list
with interrupts disabled, so its time grows with fragmentation; it is meant for
start-up allocations and for modules whose binary libraries need an allocator
(nornWIFI), not for the data path. \code{rtos\_heapStats} reports the low water mark and
the largest free block.
